%% ma: L12-B15-0006
%% lop: 12
%% chuong: 5
%% bai: 15
%% dang: 12.15.4
%% loai: TN4
%% muc: NB
%% dap_an: "A"
%% nguon: "(NBV)-ĐỀ SỐ 9-MỖI NGÀY 1 ĐỀ THI 2025.pdf (D09, MỖI NGÀY 1 ĐỀ - Đề số 9), Phần I Câu 12"
%% kien_thuc: "Bài 15 (phương trình tham số của đường thẳng, vectơ chỉ phương, hai đường thẳng song song)"
%% trang_thai: chinh-thuc
%% ghi_chu: "Đáp án gốc A đúng."
\begin{ex}
Trong không gian $Oxyz$, đường thẳng $d$ đi qua điểm $M(1;-1;3)$ và song song với đường thẳng $d_1:\dfrac{x-2}2=\dfrac{y+1}1=\dfrac{z+3}{-1}$ có phương trình là
\choice{\True $\begin{cases}x=1+2t\\y=-1+t\\z=3-t\end{cases}$}{$\begin{cases}x=1+2t\\y=-1+t\\z=3+t\end{cases}$}{$\begin{cases}x=2+t\\y=1-t\\z=-1+3t\end{cases}$}{$\begin{cases}x=1+2t\\y=1+t\\z=3-t\end{cases}$}
\loigiai{$d_1$ có vectơ chỉ phương $\overrightarrow u=(2;1;-1)$. Vì $d\parallel d_1$ nên $d$ nhận $\overrightarrow u$ làm vectơ chỉ phương và đi qua $M(1;-1;3)$, nên có phương trình tham số $x=1+2t$, $y=-1+t$, $z=3-t$.}
\end{ex}
